\section{Related Work}

\textbf{RAG systems.}
%% Definition of RAG systems
%A RAG system grounds LLM generations by retrieving relevant information from a pre-defined corpus, which may span multiple documents depending on the complexity of the user request  \cite{lewis2021retrievalaugmentedgenerationknowledgeintensivenlp, asai2020multipledocuments}. 
RAG systems extend language models with access to external knowledge from retrieved evidence \cite{lewis2021retrievalaugmentedgenerationknowledgeintensivenlp, asai2020multipledocuments, soudani2024surveyrecentadvancesconversational}. While this approach brings advantages in grounding the generation into real information \cite{askari2025selfseedingmultiintentselfinstructingllms}, it treats the query as a unitary piece of information. There are variants of RAG systems that handle complex queries, such as multi-step RAG systems which perform iterative retrieval \cite{Gu2025ChainOfRetrieval}, or feedback-based retrieval implemented with RLHF or model-estimated policies for re-generating the query until a satisfying answer can be composed by the system \cite{deng2022rlpromptoptimizingdiscretetext, rafailov2024DPO, Jun2025SearchR1}. This line of work follows earlier efforts in RL-based query reformulation \cite{Buck2018ReformulationQuery}, which have been extended to multi-modal and retrieval tasks \cite{odijkMaarten}.

\header{Query decomposition} %Complex queries can be processed by breaking them down into sub-queries for which retrieval is applied in parallel \cite{islam2024decomp}. 
Research in discourse phenomena such as exclusion \cite{excluir}, negation \cite{nevir, petcu2025comprehensivetaxonomynegationnlp}, and compositions of logic operators \cite{Zhang2024BoolQuestionsDD} shows that such formulations are difficult to encode by retrieval models \cite{krasakis2025constructingsetcompositionalnegatedrepresentations}. These cases often benefit from decomposing the original request into simpler, unitary sub-queries \cite{yaoreasoningpaths}. Query decomposition can be performed semantically using external models \cite{khot2023decomposedpromptingmodularapproach}.
However, naively splitting a complex query and retrieving documents for each sub-query does not necessarily improve results: longer contexts can degrade downstream performance \cite{shao2025reasonirtrainingretrieversreasoning}. To address this, we propose an adaptive allocation of retrieval budgets across sub-queries. %only as much evidence as needed.


\header{Research gap} We position our work at the intersection of query decomposition, using LLMs for unitary information needs, and efficient retrieval of documents by observing and modeling relevance-aware distributions for each decomposed sub-query. The closest existing work applies multi-armed bandits and active learning techniques for document selection for large-scale evaluation, and for building fair IR test collections~\cite{Li_2017, Rahman_2020, 10.1145/3269206.3271766}.